% !TEX program = xelatex
\documentclass{misisarticle}

% ============================================================
% Классификационные коды
% ============================================================

\grnti{28.23.27}
\asjc{}
\oecd{}

% ============================================================
% Название статьи
% ============================================================

\title{Условия наблюдаемости состояния БПЛА при локализации
по данным 2D LiDAR в туннельных средах}

% ============================================================
% Авторы
% ============================================================

\misisauthor
{А.\~К.\~Кузьмин}
{кафедра инженерной кибернетики НИТУ «МИСиС»}
{Москва, Россия}
{sasaluzmin360@gmail.com}

\misisauthor
{Р.\~Н.\~Садеков}
{кафедра инженерной кибернетики НИТУ «МИСиС»}
{Москва, Россия}
{r.sadekov@misis.ru}

% ============================================================
% Аннотация
% ============================================================

\abstract{%
Рассматривается задача локализации беспилотного летательного
аппарата по данным двумерного лазерного дальномера в туннельных
и других геометрически вырожденных средах.
В подобных условиях отдельные перемещения аппарата могут слабо
влиять на LiDAR-измерения, вследствие чего некоторые компоненты
состояния становятся плохо наблюдаемыми.
Целью работы является определение условий, при которых положение
и ориентация БПЛА могут быть восстановлены по последовательности
2D LiDAR-измерений.
Предлагается описывать измерения параметрической моделью среды
и исследовать чувствительность измерений к изменению состояния
аппарата с использованием матрицы Якоби.
Для количественной оценки предполагается использовать ранг
матрицы чувствительности и собственные значения соответствующей
информационной матрицы.
Проверка полученных условий планируется сначала в синтетических
туннельных средах с контролируемыми геометрическими параметрами,
а затем на данных, полученных в реальной туннельной среде.
}

\keywords{%
2D LiDAR, БПЛА, локализация, наблюдаемость,
геометрическая вырожденность, туннельная среда
}

\begin{document}

\maketitle

% ============================================================
\section{Введение}
% ============================================================

Автономная навигация беспилотных летательных аппаратов
в закрытых пространствах требует определения положения аппарата
при отсутствии сигналов глобальных навигационных спутниковых
систем. К таким средам относятся туннели, шахты, технические
полости и другие протяжённые замкнутые пространства.

Одним из источников информации о геометрии окружающей среды
является LiDAR. Измерения дальности могут использоваться для
оценивания перемещения аппарата относительно окружающих
поверхностей. Однако наличие большого количества LiDAR-измерений
само по себе не гарантирует возможность однозначного определения
состояния аппарата.

Для сред с повторяющейся геометрией возникает проблема
геометрической вырожденности. Например, в протяжённом прямом
туннеле перемещение вдоль его продольной оси может приводить
лишь к незначительному изменению измерений расстояния до стен.
При этом положение аппарата в поперечном направлении может
определяться значительно лучше.

Таким образом, информативность LiDAR-измерений зависит не только
от характеристик сенсора, но и от геометрической структуры
окружающего пространства.

Целью настоящей работы является определение условий
наблюдаемости состояния БПЛА при локализации по данным 2D LiDAR
в туннельных средах.

Для достижения поставленной цели необходимо решить следующие задачи:
\begin{itemize}
    \item сформировать математическую модель состояния БПЛА
    и LiDAR-измерений;
    \item определить связь между геометрией туннельной среды
    и чувствительностью измерений к изменению состояния;
    \item сформулировать критерий геометрической вырожденности;
    \item провести численную проверку полученных условий
    на синтетических туннельных средах;
    \item проверить применимость критерия на реальных данных.
\end{itemize}

% ============================================================
\section{Анализ существующих работ}
% ============================================================

Проблема недостаточной информативности дальномерных измерений
в геометрически вырожденных средах рассматривается в ряде работ
по LiDAR-локализации и построению карт.

В работе Zhen и Scherer~\cite{zhen2019} рассматривается локализуемость
робота в туннельной среде. Авторы представляют LiDAR-локализацию
как систему геометрических ограничений между измеренными точками
и поверхностями карты. Для оценки локализуемости исследуется
чувствительность измерений к поступательному и вращательному
движению робота. Показано, что в протяжённой туннельной среде
отдельные направления движения могут быть существенно слабее
ограничены LiDAR-измерениями, чем остальные.

В работе LION~\cite{tagliabue2020} проблема наблюдаемости
рассматривается в составе LiDAR-инерциальной системы навигации
для сред с ограниченной или отсутствующей визуальной информацией.
Оценка наблюдаемости используется для повышения устойчивости
навигационной системы в геометрически неблагоприятных условиях.

Метод X-ICP~\cite{tuna2024} развивает подход к оценке
локализуемости непосредственно в процессе регистрации облаков
точек. Информация о направлениях с различной степенью
локализуемости используется при оптимизации положения робота.

В работе Епифанова~\cite{epifanov2024} рассматриваются алгоритмы
замыкания цикла в лидарной одометрии как средство уменьшения
накопленной ошибки при построении траектории.

В работе Северина и Солдатовой~\cite{severin2026} рассматривается
построение локальной карты свободного пространства на основе
данных глубины. Работа демонстрирует использование геометрического
представления окружающего пространства для решения задач
автономной навигации мобильных робототехнических систем.

Перечисленные подходы показывают существенное влияние геометрии
окружающей среды на качество локализации. В настоящей работе
предлагается отдельно исследовать связь между параметрами
туннельной геометрии и наблюдаемостью компонент состояния БПЛА
при использовании 2D LiDAR.

% ============================================================
\section{Постановка задачи}
% ============================================================

Рассмотрим движение БПЛА в горизонтальной плоскости туннельной
среды. Предполагается, что высота полёта, а также углы крена
и тангажа стабилизируются бортовой системой управления и могут
оцениваться по другим бортовым датчикам.

Состояние, определяемое по данным 2D LiDAR, представим вектором

\begin{equation}
    \mathbf{x} =
    \begin{bmatrix}
        x & y & \psi
    \end{bmatrix}^{T},
    \label{eq:state}
\end{equation}

где $x$ и $y$ --- координаты аппарата в горизонтальной плоскости,
а $\psi$ --- угол рыскания.

Набор измерений 2D LiDAR обозначим вектором

\begin{equation}
    \mathbf{z} =
    \begin{bmatrix}
        \rho_1 & \rho_2 & \ldots & \rho_m
    \end{bmatrix}^{T},
    \label{eq:measurements}
\end{equation}

где $\rho_i$ --- дальность, измеренная $i$-м лучом LiDAR.

Геометрию окружающей среды будем задавать набором параметров
$\boldsymbol{\theta}$. В него могут входить ширина туннеля,
кривизна его продольной оси и параметры локальных геометрических
неоднородностей.

Тогда модель измерений может быть записана в общем виде

\begin{equation}
    \mathbf{z}
    =
    \mathbf{h}(\mathbf{x},\boldsymbol{\theta})
    +
    \mathbf{v},
    \label{eq:model}
\end{equation}

где $\mathbf{h}$ --- функция измерений, определяемая состоянием
аппарата и геометрией среды, а $\mathbf{v}$ --- ошибка измерений.

Требуется определить, при каких значениях
$\boldsymbol{\theta}$ изменения состояния $\mathbf{x}$
могут быть однозначно обнаружены по LiDAR-измерениям,
а при каких конфигурациях среды возникают слабонаблюдаемые
или ненаблюдаемые направления.

% ============================================================
\section{Метод исследования}
% ============================================================

Для анализа чувствительности LiDAR-измерений к изменению состояния
рассмотрим матрицу Якоби функции измерений

\begin{equation}
    \mathbf{J}
    =
    \frac{\partial \mathbf{h}}
         {\partial \mathbf{x}}.
    \label{eq:jacobian}
\end{equation}

Каждый столбец матрицы $\mathbf{J}$ характеризует изменение
LiDAR-измерений при малом изменении соответствующей компоненты
состояния.

Если существует ненулевое изменение состояния
$\Delta\mathbf{x}$, для которого

\begin{equation}
    \mathbf{J}\Delta\mathbf{x} = \mathbf{0},
    \label{eq:nullspace}
\end{equation}

то такое изменение в линейном приближении не влияет
на измерения LiDAR. Соответствующее направление состояния
является ненаблюдаемым для данной конфигурации среды.

Для состояния размерности $n$ условием локальной полной
наблюдаемости в рассматриваемой модели является

\begin{equation}
    \operatorname{rank}(\mathbf{J}) = n.
    \label{eq:rank}
\end{equation}

При наличии информации о дисперсии измерений может быть построена
информационная матрица

\begin{equation}
    \mathbf{H}
    =
    \mathbf{J}^{T}
    \mathbf{W}
    \mathbf{J},
    \label{eq:information}
\end{equation}

где $\mathbf{W}$ --- весовая матрица измерений.

Собственные значения матрицы $\mathbf{H}$ позволяют оценить
степень ограничения состояния по различным направлениям.
Малое минимальное собственное значение

\begin{equation}
    \lambda_{\min}(\mathbf{H}) \rightarrow 0
    \label{eq:eigen}
\end{equation}

соответствует появлению слабонаблюдаемого направления.

Основной эксперимент предполагается проводить на синтетических
туннельных средах. Это позволяет независимо изменять геометрические
характеристики среды и исследовать их влияние на
$\operatorname{rank}(\mathbf{J})$ и
$\lambda_{\min}(\mathbf{H})$.

После этого полученные зависимости предполагается проверить
на реальных данных LiDAR в туннельной среде.

% ============================================================
% Эксперимент и заключение пока сознательно не добавляются.
% По требованиям НИР-7 статья на данном этапе может их
% не содержать.
% ============================================================

% ============================================================
\begin{thebibliography}{9}
% ============================================================

\bibitem{zhen2019}
W. Zhen, S. Scherer.
Estimating the Localizability in Tunnel-like Environments using
LiDAR and UWB.
\textit{2019 International Conference on Robotics and Automation
(ICRA)}, pp.~4903--4908, 2019.
doi: 10.1109/ICRA.2019.8794167.

\bibitem{tagliabue2020}
A. Tagliabue, J. Tordesillas, X. Cai,
A. Santamaria-Navarro, J. P. How, L. Carlone,
A. Agha-mohammadi.
LION: Lidar-Inertial Observability-Aware Navigator
for Vision-Denied Environments.
\textit{Experimental Robotics},
pp.~380--390, 2020.
doi: 10.1007/978-3-030-71151-1\_34.

\bibitem{tuna2024}
T. Tuna, J. Nubert, Y. Nava, S. Khattak, M. Hutter.
X-ICP: Localizability-Aware LiDAR Registration
for Robust Localization in Extreme Environments.
\textit{IEEE Transactions on Robotics},
vol.~40, pp.~452--471, 2024.
doi: 10.1109/TRO.2023.3335691.

\bibitem{epifanov2024}
К.~А.~Епифанов.
Исследование алгоритмов замыкания цикла
в лидарной одометрии.
\textit{Искусственный интеллект в промышленных,
коммерческих, медицинских и финансовых приложениях:
сборник статей научно-технического семинара студентов
кафедры «Инженерной кибернетики»}.
Москва: НИТУ МИСИС, 2024, с.~57--61.

\bibitem{severin2026}
Я.~А.~Северин, Н.~Е.~Солдатова.
Построение локальной карты свободного пространства
в помещении на основе RGB-D данных.
\textit{Искусственный интеллект в промышленных,
коммерческих, медицинских и финансовых приложениях:
сборник статей научно-технологического семинара
кафедры «Инженерной кибернетики»}.
Вып.~5.
Москва: НИТУ МИСИС, 2026, с.~105--112.

% TODO:
% Добавить одну публикацию преподавателя.
% Это отдельно требуется в таблице НИР-7.

\end{thebibliography}

\end{document}
